\documentclass[11pt,a4paper]{article}

\usepackage[margin=2.5cm]{geometry}
\usepackage{amsmath,amssymb,amsthm}
\usepackage{booktabs}
\usepackage{array}
\usepackage{hyperref}
\usepackage[most]{tcolorbox}
\usepackage{pdflscape}

\theoremstyle{plain}
\newtheorem{theorem}{Theorem}[section]
\newtheorem{proposition}[theorem]{Proposition}
\newtheorem{corollary}[theorem]{Corollary}
\newtheorem{lemma}[theorem]{Lemma}
\newtheorem{conjecture}[theorem]{Conjecture}
\theoremstyle{definition}
\newtheorem{definition}[theorem]{Definition}
\newtheorem{observation}[theorem]{Observation}
\newtheorem{remark}[theorem]{Remark}

\newcommand{\R}{\mathbb{R}}
\newcommand{\Z}{\mathbb{Z}}
\newcommand{\N}{\mathbb{N}}
\newcommand{\C}{\mathbb{C}}

\title{\textbf{Matching Matrices:\\
From Feynman Diagrams to Qubits}}
\author{Robin Langer and Claude Code}
\date{August 2026}

\begin{document}
\maketitle

\begin{abstract}
We study the change-of-basis matrices between orthogonal polynomials
and their associated orthonormal functions --- the eigenfunctions of
a symmetry operator --- in both the continuous and discrete settings.

In the continuous case, the Hermite functions (eigenfunctions of the
Fourier transform) projected onto the Hermite polynomial basis yield
an infinite matrix whose entries, after rescaling by $3^{(m+k)/2}$,
are integers $D_{mk}$ that count signed bipartite matchings ---
equivalently, Feynman diagrams. The generating function is
$\sqrt{2\pi/3}\;\exp(-(s^2 - 4su + u^2)/3)$.

In the discrete case, the Krawtchouk polynomials (eigenvectors of
the Walsh--Hadamard transform on $n$ qubits) evaluated against a
binomial weight at parameter $q$ yield a finite matrix $I(m,k,n,q)$
whose generating function is $[1 + \alpha(s+u) + su]^n$, where
$\alpha = 1 - 2q$. For integer $\alpha$, every entry is
automatically an integer --- a weighted count of multinomial
microstates.

The two theories are connected by a scaling limit: as $n \to \infty$,
the binomial converges to the Gaussian, Krawtchouk converges to
Hermite, and multinomial coefficients converge to matchings. The
passage from the discrete matrix to the continuous matrix is the
passage from an $n$-qubit spin system to quantum field theory.
\end{abstract}

%% ===================================================================
\section{The Construction}
%% ===================================================================

\begin{tcolorbox}[breakable, colback=black!5, colframe=black!5,
  boxrule=0pt, sharp corners]
Both cases follow the same pattern, in three steps.

\textbf{Step 1: the polynomials.} Fix a weight $W \ge 0$ on a set
$X$. It defines an inner product on polynomials,
$\langle f, g \rangle_W = \int_X f\, g\; W$.
Gram--Schmidt produces the orthogonal polynomials $P_n$:
$\langle P_m, P_n \rangle_W = h_n \delta_{mn}$.

\textbf{Step 2: the functions.} Multiplying by $\sqrt{W}$ turns
them into the \emph{associated functions}
$\Phi_n = \widehat P_n\, \sqrt{W}$,
which are orthonormal for the plain inner product:
$\int_X \Phi_m\, \Phi_n = \delta_{mn}$.

\textbf{Step 3: the coefficients.} Project $\Phi_m$ back onto
$\{P_k\}$ in $L^2(W)$:
\[
  b_{mk} = \langle \Phi_m, \widehat P_k \rangle_W
  = \int_X \widehat P_m\, \widehat P_k\; W^{3/2}.
\]

The factor $W^{3/2}$ arises from $\sqrt{W}$ (from $\Phi_m$)
times $W$ (from the inner product). The matrix $(b_{mk})$
is our object of study. The question: is there a rescaling
that makes it integer, and if so, what do the integers count?
\end{tcolorbox}

%% ===================================================================
%% PART I: CONTINUOUS
%% ===================================================================

\part*{Part I: The Continuous Case}
\addcontentsline{toc}{part}{Part I: The Continuous Case}

%% ===================================================================
\section{Hermite Polynomials and Hermite Functions}
%% ===================================================================

Let $H_n(t)$ denote the physicist's Hermite polynomials, defined by
the recurrence
\[
  H_{n+1}(t) = 2t\, H_n(t) - 2n\, H_{n-1}(t),
  \qquad H_0(t) = 1,\quad H_1(t) = 2t,
\]
and satisfying the orthogonality relation
\begin{equation}\label{eq:hermite-orth}
  \int_{-\infty}^{\infty} H_m(t)\, H_k(t)\, e^{-t^2}\, dt
  = 2^k\, k!\, \sqrt{\pi}\; \delta_{mk}.
\end{equation}
Here $W = e^{-t^2}$ on $\R$, with $h_n = 2^n n! \sqrt{\pi}$.
The \emph{Hermite functions} are
\begin{equation}\label{eq:hermite-fn}
  \psi_m(t) = c_m\, H_m(t)\, e^{-t^2/2},
  \qquad c_m = (2^m\, m!\, \sqrt{\pi})^{-1/2}.
\end{equation}
These form an orthonormal basis for $L^2(\R)$ and are eigenfunctions
of the Fourier transform with eigenvalues $(-i)^m$.

\begin{tcolorbox}[colback=yellow!10, colframe=yellow!10, boxrule=0pt,
  sharp corners]
The Fourier transform $\widehat{f}(\xi) = \int f(t)\, e^{-2\pi i
t \xi}\, dt$ is an operator of order four: $\widehat{\widehat{
\widehat{\widehat{f}}}} = f$. Its eigenvalues must therefore be
fourth roots of unity: $\{1, -i, -1, i\}$. The Hermite functions
are the eigenfunctions, partitioned into four eigenspaces according
to $m \bmod 4$. The fact that $\psi_m$ is a Fourier eigenfunction is
the statement that the quantum harmonic oscillator is invariant under
the exchange of position and momentum: $x \leftrightarrow p$.
\end{tcolorbox}

The $L^2$-projection of $\psi_m$ onto the Hermite polynomial basis
gives coefficients
\begin{equation}\label{eq:amk}
  a_{mk}
  = \frac{c_m}{2^k\, k!\,\sqrt{\pi}}
    \int_{-\infty}^{\infty} H_m(t)\, H_k(t)\, e^{-3t^2/2}\, dt.
\end{equation}
The factor $e^{-3t^2/2}$ arises from $e^{-t^2/2}$ (from $\psi_m$)
times $e^{-t^2}$ (from the weight).

%% ===================================================================
\section{The Hermite Generating Function}
%% ===================================================================

\begin{tcolorbox}[colback=black!5, colframe=black!5, boxrule=0pt,
  sharp corners]
\begin{proposition}\label{prop:hermite-gf}
The integrals $I_{mk} = \int_{-\infty}^{\infty} H_m(t)\, H_k(t)\,
e^{-3t^2/2}\, dt$ have the generating function
\[
  \sum_{m,k \ge 0} \frac{s^m\, u^k}{m!\, k!}\; I_{mk}
  = \sqrt{\frac{2\pi}{3}}\;
    \exp\!\Bigl(-\frac{s^2 - 4su + u^2}{3}\Bigr).
\]
\end{proposition}
\end{tcolorbox}

\begin{proof}
Using the generating function for Hermite polynomials,
$\sum_{n \ge 0} H_n(t)\, z^n / n! = e^{2tz - z^2}$:
\begin{align*}
  \sum_{m,k} \frac{s^m u^k}{m!\,k!}\, I_{mk}
  &= \int_{-\infty}^{\infty}
     e^{2ts - s^2}\, e^{2tu - u^2}\, e^{-3t^2/2}\, dt \\
  &= \int_{-\infty}^{\infty}
     \exp\!\bigl(-\tfrac{3}{2}t^2 + 2(s+u)t - s^2 - u^2\bigr)\, dt \\
  &= \sqrt{\frac{2\pi}{3}}\;
     \exp\!\Bigl(\frac{2(s+u)^2}{3} - s^2 - u^2\Bigr)
  = \sqrt{\frac{2\pi}{3}}\;
     \exp\!\Bigl(-\frac{s^2 - 4su + u^2}{3}\Bigr).
     \qedhere
\end{align*}
\end{proof}

\begin{tcolorbox}[colback=black!5, colframe=black!5, boxrule=0pt,
  sharp corners]
\begin{corollary}[Parity]\label{cor:parity}
$I_{mk} = 0$ whenever $m + k$ is odd.
\end{corollary}
\end{tcolorbox}

\begin{proof}
The exponent $-(s^2 - 4su + u^2)/3$ contains only terms of even
total degree in $(s,u)$.
\end{proof}

%% ===================================================================
\section{The Hermite Integer Matrix}
%% ===================================================================

\begin{tcolorbox}[colback=black!5, colframe=black!5, boxrule=0pt,
  sharp corners]
\begin{theorem}[Hermite integrality]\label{thm:hermite-int}
Define
$D_{mk} = 3^{(m+k)/2} \cdot I_{mk} / \sqrt{2\pi/3}$.
Then $D_{mk} \in \Z$ for all $m, k \ge 0$ with $m + k$ even
(and zero when $m + k$ is odd). Explicitly, with
$b = (m-a)/2$, $c = (k-a)/2$,
\[
  D_{mk} = \sum_{a}\; 4^a\, (-2)^{b+c}\, a!\;
  \frac{m!}{a!\, b!\, 2^b}\; \frac{k!}{a!\, c!\, 2^c}\,,
\]
where each factor $n!/(a!\, b!\, 2^b)$ counts the partial matchings
of $[n]$ with $a$ unmatched points and $b$ pairs.
\end{theorem}
\end{tcolorbox}

\begin{proof}
Write
$\exp(-(s^2 - 4su + u^2)/3) = \exp(-s^2/3)\, \exp(4su/3)\,
\exp(-u^2/3)$.
The coefficient of $s^m u^k$ is a sum over triples $(a, b, c)$ with
$m = 2b + a$ and $k = a + 2c$:
\[
  [s^m u^k] = \sum_{a,b,c}
  \frac{(-1/3)^b}{b!}\;\frac{(4/3)^a}{a!}\;\frac{(-1/3)^c}{c!}.
\]
Multiplying by $m!\, k!\, 3^{(m+k)/2}$ absorbs all powers of~$3$
and leaves integer factors.
\end{proof}

\begin{tcolorbox}[colback=blue!5, colframe=blue!5, boxrule=0pt,
  sharp corners]
\begin{conjecture}[Hermite matching conjecture]\label{conj:hermite}
$D_{mk}$ counts signed bipartite matchings. Colour $m$ points red
and $k$ points blue; then
\[
  D_{mk} = \sum_{\substack{\text{perfect matchings}\\
  \text{of the } m+k \text{ points}}}
  4^{\#\{\text{red--blue pairs}\}} \; (-2)^{\#\{\text{same-colour pairs}\}}.
\]
Equivalently,
\[
  D_{mk} = \sum_{\substack{c \ge 0 \\ m - c,\, k - c \text{ even}}}
  \binom{m}{c}\binom{k}{c}\, c!\;
  (m{-}c{-}1)!!\; (k{-}c{-}1)!!\; 4^{c}\, (-2)^{(m+k-2c)/2}.
\]
Verified for all $m, k \le 20$ (441 entries, zero discrepancies).
The matching interpretation is consistent with Wick's theorem applied
to $\mathbb{E}[H_m(X)\, H_k(X)]$ for $X \sim N(0, 1/3)$.
\end{conjecture}
\end{tcolorbox}

%% -------------------------------------------------------------------
\subsection{The diagonal}

\begin{tcolorbox}[colback=black!5, colframe=black!5, boxrule=0pt,
  sharp corners]
\begin{theorem}[Hermite diagonal]\label{thm:hermite-diag}
\begin{equation}\label{eq:hermite-diag}
  \frac{D_{nn}}{n!} = \sum_{j=0}^{\lfloor n/2 \rfloor}
  \binom{n}{2j}\binom{2j}{j}\, 4^{n-2j}
  = 1,\, 4,\, 18,\, 88,\, 454,\, 2424,\, \ldots
  \quad (\text{OEIS A081671}),
\end{equation}
with ordinary generating function
$1/\sqrt{(1-2x)(1-6x)}$
and exponential generating function $e^{4x}\, I_0(2x)$, where $I_0$
is the modified Bessel function of the first kind.
\end{theorem}
\end{tcolorbox}

\begin{proof}
Setting $m = k = n$ in the matching formula and simplifying with
$(2j-1)!! = (2j)!/(2^j\, j!)$ yields the stated sum. The OGF
follows from $\sum_n a_n x^n$ using the Cauchy product of
$1/\sqrt{1-4z^2}$ with $1/(1-4x)$, giving
$1/\sqrt{(1-4x)^2 - 4x^2} = 1/\sqrt{1 - 8x + 12x^2}$.
\end{proof}

%% -------------------------------------------------------------------
\subsection{Edge entries}

\begin{tcolorbox}[colback=black!5, colframe=black!5, boxrule=0pt,
  sharp corners]
\begin{proposition}[Hermite edge]\label{prop:hermite-edge}
$D_{0,2n} = (-1)^n\, (2n)!/n! = (-1)^n\, 2^n\, (2n{-}1)!!$.
\end{proposition}
\end{tcolorbox}

%% -------------------------------------------------------------------
\subsection{Reverse direction}

The \emph{reverse} expansion $\widehat H_m = \sum_k \widetilde
b_{mk}\, \psi_k$ has one power of $\sqrt{W}$ in place of three:
$\widetilde b_{mk} = \int \widehat H_m\, \widehat H_k\,
e^{-t^2/2}\, dt$. The same rescaling without the power of~$3$ gives

\begin{tcolorbox}[colback=black!5, colframe=black!5, boxrule=0pt,
  sharp corners]
\begin{theorem}[Hermite reverse]\label{thm:hermite-rev}
$S_{mk} = (-1)^{(k-m)/2}\, D_{mk}$ --- the forward matrix with the
signs stripped:
\[
  S_{mk} = \sum_{\substack{\text{perfect matchings}\\
  \text{of the } m+k \text{ points}}}
  4^{\#\{\text{red--blue}\}} \; 2^{\#\{\text{same-colour}\}}.
\]
Every entry is non-negative.
\end{theorem}
\end{tcolorbox}

%% -------------------------------------------------------------------
\subsection{Computed values}

\begin{tcolorbox}[colback=black!5, colframe=black!5, boxrule=0pt,
  sharp corners]
\begin{center}
$D_{mk}$ for even $m$ and $k$, $0 \le m,k \le 14$

\medskip
\scriptsize
\setlength{\tabcolsep}{3pt}
\begin{tabular}{r|rrrrrr}
$m \!\backslash\! k$ & $0$ & $2$ & $4$ & $6$ & $8$ & $10$ \\
\hline
$0$ & $1$ & $-2$ & $12$ & $-120$ & $1680$ & $-30240$ \\
$2$ & $-2$ & $36$ & $-408$ & $6000$ & $-110880$ & $2479680$ \\
$4$ & $12$ & $-408$ & $10896$ & $-254880$ & $6471360$ & $-184222080$ \\
$6$ & $-120$ & $6000$ & $-254880$ & $9529920$ & $-339534720$ & $12515731200$ \\
$8$ & $1680$ & $-110880$ & $6471360$ & $-339534720$ & $16590147840$ \\
$10$ & $-30240$ & $2479680$ & $-184222080$ & $12515731200$ \\
\end{tabular}
\end{center}
\end{tcolorbox}

%% ===================================================================
%% PART II: DISCRETE
%% ===================================================================

\part*{Part II: The Discrete Case}
\addcontentsline{toc}{part}{Part II: The Discrete Case}

%% ===================================================================
\section{Krawtchouk Polynomials}
%% ===================================================================

The \emph{Krawtchouk polynomials} $K_j(x; n)$ for $p = \tfrac12$
are defined on $\{0, 1, \ldots, n\}$ by
\begin{equation}\label{eq:kraw-def}
  K_j(x; n) = \sum_{r=0}^{j} (-1)^r \binom{x}{r}
  \binom{n-x}{j-r},
\end{equation}
with generating function
\begin{equation}\label{eq:kraw-gf}
  \sum_{j=0}^{n} K_j(x; n)\, z^j = (1-z)^x\, (1+z)^{n-x},
\end{equation}
and orthogonality
\begin{equation}\label{eq:kraw-orth}
  \sum_{x=0}^{n} K_j(x; n)\, K_k(x; n)\, \frac{\binom{n}{x}}{2^n}
  = \binom{n}{k}\, \delta_{jk}.
\end{equation}

\begin{tcolorbox}[colback=yellow!10, colframe=yellow!10, boxrule=0pt,
  sharp corners]
The Krawtchouk polynomials are the discrete Hermite polynomials: they
are orthogonal with respect to the binomial distribution (the discrete
Gaussian), and they converge to Hermite polynomials in the
$n \to \infty$ limit. The matrix $(K_j(k; n))_{j,k=0}^n$ is, up to
normalisation, the Walsh--Hadamard transform --- the Fourier transform
of $(\mathbb{Z}/2\mathbb{Z})^n$, i.e.\ the Hadamard gate applied to
$n$ qubits.
\end{tcolorbox}

%% ===================================================================
\section{The Krawtchouk Matching Matrix}
%% ===================================================================

\begin{tcolorbox}[colback=green!8, colframe=green!8, boxrule=0pt,
  sharp corners]
\begin{definition}[Krawtchouk matching matrix]\label{def:kraw}
Fix $n \ge 1$ and a parameter $q$. The \emph{Krawtchouk matching
matrix} is
\begin{equation}\label{eq:Imk}
  I(m, k, n, q)
  = \sum_{x=0}^{n} K_m(x; n)\, K_k(x; n)\, \binom{n}{x}\,
  q^x\, (1-q)^{n-x}
  = \mathbb{E}\bigl[K_m(X; n)\, K_k(X; n)\bigr]
\end{equation}
for $X \sim \mathrm{Bin}(n, q)$.
\end{definition}
\end{tcolorbox}

At $q = \tfrac12$ this reduces to the orthogonality
relation~\eqref{eq:kraw-orth}: the matrix is diagonal. For
$q \ne \tfrac12$ the off-diagonal entries are nonzero.

%% ===================================================================
\section{The Krawtchouk Generating Function}
%% ===================================================================

\begin{tcolorbox}[colback=black!5, colframe=black!5, boxrule=0pt,
  sharp corners]
\begin{theorem}[Bivariate generating function]\label{thm:kraw-gf}
Set $\alpha = 1 - 2q$. Then
\begin{equation}\label{eq:kraw-gf-main}
  F(s, u) = \sum_{m,k=0}^{n} I(m, k, n, q)\, s^m\, u^k
  = \bigl[\,1 + \alpha(s + u) + su\,\bigr]^n.
\end{equation}
\end{theorem}
\end{tcolorbox}

\begin{proof}
Using the Krawtchouk generating function~\eqref{eq:kraw-gf}:
\begin{align*}
  F(s,u)
  &= \sum_{x=0}^{n} \binom{n}{x}\, q^x\, (1{-}q)^{n-x}\,
     (1{-}s)^x(1{+}s)^{n-x}\, (1{-}u)^x(1{+}u)^{n-x} \\
  &= \bigl[\, q(1{-}s)(1{-}u) + (1{-}q)(1{+}s)(1{+}u) \,\bigr]^n
  = \bigl[\, 1 + \alpha(s{+}u) + su \,\bigr]^n.
  \qedhere
\end{align*}
\end{proof}

\begin{tcolorbox}[breakable, colback=yellow!10, colframe=yellow!10,
  boxrule=0pt, sharp corners]
Compare the two generating functions:

\medskip
\begin{center}
\renewcommand{\arraystretch}{1.5}
\begin{tabular}{lcc}
  \toprule
  & Continuous (Hermite) & Discrete (Krawtchouk) \\
  \midrule
  GF & $\sqrt{\tfrac{2\pi}{3}}\,
    \exp\!\bigl(-\frac{s^2 - 4su + u^2}{3}\bigr)$
     & $[1 + \alpha(s{+}u) + su]^n$ \\[6pt]
  Shape & exp of quadratic form & $n$-th power \\
  Parameter & $\sigma^2 = 1/3$ (fixed) & $\alpha \in \Z$ (free) \\
  \bottomrule
\end{tabular}
\end{center}

\medskip
The continuous generating function is an \emph{exponential} of a
quadratic form; the discrete is an $n$-th \emph{power} of an
affine--bilinear form. They are connected by the usual
$\exp = \lim (1 + x/n)^n$: scaling $s$, $u$, $\alpha$ by
$1/\sqrt{n}$ and sending $n \to \infty$, the power converges to the
exponential.
\end{tcolorbox}

%% ===================================================================
\section{Coefficient Extraction}
%% ===================================================================

\begin{tcolorbox}[colback=black!5, colframe=black!5, boxrule=0pt,
  sharp corners]
\begin{theorem}[Multinomial formula]\label{thm:kraw-coeff}
\begin{equation}\label{eq:kraw-coeff}
  I(m, k, n, q)
  = \sum_{c=\max(0,m+k-n)}^{\min(m,k)}
  \frac{n!}{(n{-}m{-}k{+}c)!\;(m{-}c)!\;(k{-}c)!\;c!}
  \;\alpha^{\,m+k-2c}.
\end{equation}
\end{theorem}
\end{tcolorbox}

\begin{proof}
Expand $[1 + \alpha(s{+}u) + su]^n$ by the multinomial theorem with
three terms. Pick $a$ copies of~$1$, $b$ copies of~$\alpha(s{+}u)$,
and $c$ copies of~$su$. Then $a + b + c = n$, and the monomial
$s^m u^k$ requires $b = m + k - 2c$ and $a = n - m - k + c$.
\end{proof}

\begin{tcolorbox}[colback=black!5, colframe=black!5, boxrule=0pt,
  sharp corners]
\begin{corollary}[Automatic integrality]\label{cor:kraw-int}
When $\alpha \in \mathbb{Z}$, every entry
$I(m, k, n, q) \in \mathbb{Z}$. No rescaling is required.
\end{corollary}
\end{tcolorbox}

%% ===================================================================
\section{Properties}
%% ===================================================================

\begin{tcolorbox}[colback=black!5, colframe=black!5, boxrule=0pt,
  sharp corners]
\begin{proposition}[Symmetries]\label{prop:kraw-sym}
\mbox{}
\begin{enumerate}
\item \textbf{Transpose:} $I(m,k) = I(k,m)$.
\item \textbf{Complement:}
  $I(m,k,n,q) = \alpha^n\, I(n{-}m,\, n{-}k,\, n,\, q)$.
\item \textbf{Parameter reflection:}
  $I(m,k,n,1{-}q) = (-1)^{m+k}\, I(m,k,n,q)$.
\item \textbf{No parity vanishing:}
  entries with $m+k$ odd are generically nonzero (unlike Hermite).
\end{enumerate}
\end{proposition}
\end{tcolorbox}

\begin{tcolorbox}[colback=black!5, colframe=black!5, boxrule=0pt,
  sharp corners]
\begin{proposition}[Edge entries]\label{prop:kraw-edge}
$I(0, k, n, q) = \binom{n}{k}\, \alpha^k$.
\end{proposition}
\end{tcolorbox}

\begin{tcolorbox}[colback=black!5, colframe=black!5, boxrule=0pt,
  sharp corners]
\begin{proposition}[Special values]\label{prop:kraw-special}
\mbox{}
\begin{enumerate}
\item $\alpha = 0$ ($q = \tfrac12$):
  $F = (1+su)^n$, $I(m,k) = \binom{n}{k}\,\delta_{mk}$
  (orthogonality).
\item $\alpha = 1$ ($q = 0$):
  $F = [(1{+}s)(1{+}u)]^n$, $I(m,k) = \binom{n}{m}\binom{n}{k}$.
\item $\alpha = -1$ ($q = 1$):
  $F = [(1{-}s)(1{-}u)]^n$,
  $I(m,k) = (-1)^{m+k}\binom{n}{m}\binom{n}{k}$.
\end{enumerate}
\end{proposition}
\end{tcolorbox}

\begin{tcolorbox}[colback=black!5, colframe=black!5, boxrule=0pt,
  sharp corners]
\begin{proposition}[Rational rescaling]\label{prop:kraw-rational}
If $\alpha = p/d$ with $p, d \in \mathbb{Z}$, then
$d^{m+k}\, I(m, k, n, q) \in \mathbb{Z}$.
\end{proposition}
\end{tcolorbox}

%% -------------------------------------------------------------------
\subsection{Computed values}

\begin{tcolorbox}[colback=black!5, colframe=black!5, boxrule=0pt,
  sharp corners]
\begin{center}
$I(m,k,n,q)$ for $\alpha = 2$, $n = 6$

\medskip
\begin{tabular}{r|rrrrrrr}
$m \!\backslash\! k$ & $0$ & $1$ & $2$ & $3$ & $4$ & $5$ & $6$ \\
\hline
$0$ & $1$ & $12$ & $60$ & $160$ & $240$ & $192$ & $64$ \\
$1$ & $12$ & $126$ & $540$ & $1200$ & $1440$ & $864$ & $192$ \\
$2$ & $60$ & $540$ & $1935$ & $3480$ & $3240$ & $1440$ & $240$ \\
$3$ & $160$ & $1200$ & $3480$ & $4900$ & $3480$ & $1200$ & $160$ \\
$4$ & $240$ & $1440$ & $3240$ & $3480$ & $1935$ & $540$ & $60$ \\
$5$ & $192$ & $864$ & $1440$ & $1200$ & $540$ & $126$ & $12$ \\
$6$ & $64$ & $192$ & $240$ & $160$ & $60$ & $12$ & $1$ \\
\end{tabular}
\end{center}
\end{tcolorbox}

\begin{tcolorbox}[colback=black!5, colframe=black!5, boxrule=0pt,
  sharp corners]
\begin{center}
$I(m,k,n,q)$ for $\alpha = 2$, $n = 8$

\medskip
\small
\begin{tabular}{r|rrrrrrrrr}
$m \!\backslash\! k$ & $0$ & $1$ & $2$ & $3$ & $4$ & $5$ & $6$ & $7$ & $8$ \\
\hline
$0$ & $1$ & $16$ & $112$ & $448$ & $1120$ & $1792$ & $1792$ & $1024$ & $256$ \\
$1$ & $16$ & $232$ & $1456$ & $5152$ & $11200$ & $15232$ & $12544$ & $5632$ & $1024$ \\
$2$ & $112$ & $1456$ & $8092$ & $24976$ & $46480$ & $52864$ & $35392$ & $12544$ & $1792$ \\
$3$ & $448$ & $5152$ & $24976$ & $66136$ & $103600$ & $97216$ & $52864$ & $15232$ & $1792$ \\
$4$ & $1120$ & $11200$ & $46480$ & $103600$ & $134470$ & $103600$ & $46480$ & $11200$ & $1120$ \\
$5$ & $1792$ & $15232$ & $52864$ & $97216$ & $103600$ & $66136$ & $24976$ & $5152$ & $448$ \\
$6$ & $1792$ & $12544$ & $35392$ & $52864$ & $46480$ & $24976$ & $8092$ & $1456$ & $112$ \\
$7$ & $1024$ & $5632$ & $12544$ & $15232$ & $11200$ & $5152$ & $1456$ & $232$ & $16$ \\
$8$ & $256$ & $1024$ & $1792$ & $1792$ & $1120$ & $448$ & $112$ & $16$ & $1$ \\
\end{tabular}
\end{center}
\end{tcolorbox}

\begin{tcolorbox}[colback=black!5, colframe=black!5, boxrule=0pt,
  sharp corners]
\begin{center}
$I(m,k,n,q)$ for $\alpha = 3$, $n = 6$

\medskip
\begin{tabular}{r|rrrrrrr}
$m \!\backslash\! k$ & $0$ & $1$ & $2$ & $3$ & $4$ & $5$ & $6$ \\
\hline
$0$ & $1$ & $18$ & $135$ & $540$ & $1215$ & $1458$ & $729$ \\
$1$ & $18$ & $276$ & $1710$ & $5400$ & $8910$ & $6804$ & $1458$ \\
$2$ & $135$ & $1710$ & $8385$ & $19620$ & $21465$ & $8910$ & $1215$ \\
$3$ & $540$ & $5400$ & $19620$ & $30800$ & $19620$ & $5400$ & $540$ \\
$4$ & $1215$ & $8910$ & $21465$ & $19620$ & $8385$ & $1710$ & $135$ \\
$5$ & $1458$ & $6804$ & $8910$ & $5400$ & $1710$ & $276$ & $18$ \\
$6$ & $729$ & $1458$ & $1215$ & $540$ & $135$ & $18$ & $1$ \\
\end{tabular}
\end{center}
\end{tcolorbox}

%% ===================================================================
%% PART III: COMPARISON AND PHYSICS
%% ===================================================================

\part*{Part III: Comparison and Physics}
\addcontentsline{toc}{part}{Part III: Comparison and Physics}

%% ===================================================================
\section{Structural Comparison}
%% ===================================================================

\begin{tcolorbox}[breakable, colback=yellow!10, colframe=yellow!10,
  boxrule=0pt, sharp corners]
\begin{center}
\renewcommand{\arraystretch}{1.5}
\begin{tabular}{lcc}
  \toprule
  Feature & Continuous (Hermite) & Discrete (Krawtchouk) \\
  \midrule
  Weight & $e^{-t^2}$ on $\R$ & $\binom{n}{x}/2^n$ on $\{0,\ldots,n\}$ \\
  Polynomials & $H_n(t)$ & $K_j(x;n)$ \\
  Functions & $\psi_n = c_n H_n e^{-t^2/2}$ & (absorbed into weight) \\
  Transform & Fourier on $\R$ & Walsh--Hadamard on $\{0,1\}^n$ \\
  Eigenvalues & $(-i)^n$ & $(-1)^j$ \\[4pt]
  GF &
    $\sqrt{\tfrac{2\pi}{3}}\,e^{-(s^2-4su+u^2)/3}$
    & $[1 + \alpha(s{+}u) + su]^n$ \\[4pt]
  Integrality & rescale by $3^{(m+k)/2}$ & automatic for $\alpha \in \Z$ \\
  Parity & $D_{mk}=0$ for $m{+}k$ odd & no vanishing \\
  Edge & $(-1)^n(2n)!/n!$ & $\binom{n}{k}\alpha^k$ \\
  At $\alpha=1$ & --- & $\binom{n}{m}\binom{n}{k}$ \\
  Symmetry & $(m,k) \leftrightarrow (k,m)$
    & transpose $+$ complement \\
  Combinatorics & signed matchings & multinomial coefficients \\
  Physics & Feynman diagrams & qubit microstates \\
  \bottomrule
\end{tabular}
\end{center}
\end{tcolorbox}

%% ===================================================================
\section{Quantum Mechanics}
%% ===================================================================

\begin{tcolorbox}[breakable, colback=yellow!10, colframe=yellow!10,
  boxrule=0pt, sharp corners]
The matching matrices are not merely combinatorial curiosities. They
sit at the centre of quantum mechanics, and the passage from the
discrete matrix to the continuous matrix is the passage from finite
spin systems to quantum field theory.
\end{tcolorbox}

%% -------------------------------------------------------------------
\subsection{The harmonic oscillator}

The Hermite functions $\psi_n$ are the energy eigenstates of the
\textbf{quantum harmonic oscillator}, the single most important
quantum system. Its Hamiltonian is $H = p^2 + x^2$ (up to
constants), with eigenvalues $E_n = (n + \tfrac12)\hbar\omega$.
Every quadratic quantum system --- photons, phonons, the free
electromagnetic field --- is a harmonic oscillator.

\begin{tcolorbox}[colback=black!5, colframe=black!5, boxrule=0pt,
  sharp corners]
\begin{observation}[Physical meaning of $D_{mk}$]\label{obs:Dmk}
The Hermite matching integer $D_{mk}$ is proportional to the overlap
between the $m$-th energy eigenstate ($\psi_m$) and the $k$-th
polynomial observable ($H_k$) in $L^2(e^{-t^2})$. Physically: if a
particle is in the $m$-th energy level, $D_{mk}$ measures the $k$-th
polynomial moment of its position, weighted by the Gaussian ground
state.
\end{observation}
\end{tcolorbox}

%% -------------------------------------------------------------------
\subsection{Matchings are Feynman diagrams}

\begin{tcolorbox}[breakable, colback=yellow!10, colframe=yellow!10,
  boxrule=0pt, sharp corners]
\textbf{Wick's theorem} states that the expectation of a product of
Gaussian random variables decomposes into a sum over all pairings:
\[
  \mathbb{E}[X_1 X_2 \cdots X_{2n}]
  = \sum_{\text{matchings}} \prod_{\text{pairs } (i,j)}
    \mathbb{E}[X_i X_j].
\]
Each matching is a \textbf{Feynman diagram}: a way of contracting
field operators into propagators. In free quantum field theory,
every scattering amplitude is a sum over such diagrams.

The Hermite matching integer $D_{mk}$ --- the sum over bipartite
matchings with weights $+4$ (cross-pairs) and $-2$ (self-pairs) ---
\emph{is} Wick's theorem applied to $H_m(X)\, H_k(X)$ at variance
$\sigma^2 = 1/3$:
\begin{itemize}
\item A \textbf{cross-pair} connects a slot from $H_m$ to one from
  $H_k$: an inter-level transition, weighted $+4$.
\item A \textbf{self-pair} connects two slots within the same
  polynomial: an intra-level contraction, weighted $-2$.
\end{itemize}
The Hermite matching matrix is a \textbf{table of Feynman diagram
counts}, indexed by energy levels.
\end{tcolorbox}

%% -------------------------------------------------------------------
\subsection{Qubits and spin}

\begin{tcolorbox}[breakable, colback=yellow!10, colframe=yellow!10,
  boxrule=0pt, sharp corners]
The Krawtchouk matching matrix has its own quantum mechanical
meaning.

Krawtchouk polynomials are the matrix elements of $SU(2)$
representations. A system of $n$ spin-$\tfrac12$ particles ($n$
qubits) has a Hilbert space of dimension $2^n$. The total spin
operator has eigenspaces labelled by $x \in \{0,\ldots,n\}$ (the
number of spin-ups), and $K_j(x;n)$ is the $j$-th matrix element.

The Walsh--Hadamard transform $(K_j(k;n))$ is the Hadamard gate
applied to $n$ qubits --- the first gate in Deutsch--Jozsa, Grover,
and Shor.
\end{tcolorbox}

\begin{tcolorbox}[colback=black!5, colframe=black!5, boxrule=0pt,
  sharp corners]
\begin{observation}[Physical meaning of $I(m,k,n,q)$]
\label{obs:Imk}
$I(m,k,n,q)$ measures the correlation between the $m$-th and $k$-th
Hadamard-basis observables when the $n$ qubits are in a biased
product state with probability $q$ of spin-up.

At $q = \tfrac12$ ($\alpha = 0$): the state is $|{+}\rangle^{\otimes
n}$. The observables are decorrelated (orthogonality).

At $q \ne \tfrac12$: the observables become correlated. The
parameter $\alpha$ measures the strength of quantum correlation.
\end{observation}
\end{tcolorbox}

%% -------------------------------------------------------------------
\subsection{The generating function as a partition function}

\begin{tcolorbox}[breakable, colback=yellow!10, colframe=yellow!10,
  boxrule=0pt, sharp corners]
The generating function $[1 + \alpha(s{+}u) + su]^n$ is a partition
function. Each of the $n$ qubits independently contributes one of
three terms:

\begin{center}
\begin{tabular}{lcl}
  \toprule
  Term & Weight & Meaning \\
  \midrule
  $1$ & $1$ & qubit idle \\
  $\alpha s$ or $\alpha u$ & $\alpha$ & qubit contributes to one
    observable \\
  $su$ & $1$ & qubit contributes to both \\
  \bottomrule
\end{tabular}
\end{center}

\medskip
This is a lattice model: $n$ independent sites, each in one of three
effective states, with a global constraint (total $s$-degree $= m$,
total $u$-degree $= k$).
\end{tcolorbox}

%% -------------------------------------------------------------------
\subsection{Why integers?}

\begin{tcolorbox}[colback=black!5, colframe=black!5, boxrule=0pt,
  sharp corners]
\begin{observation}[Integrality as discreteness of process space]
\label{obs:integers}
The integers count the number of topologically distinct quantum
processes.

\begin{itemize}
\item \textbf{Continuous:} $D_{mk}$ counts signed Feynman diagrams
  --- the distinct ways energy levels $m$ and $k$ can exchange
  virtual quanta through Gaussian contractions. The integrality says
  that the space of quantum processes has a discrete structure,
  even though the underlying physics is continuous.
\item \textbf{Discrete:} $I(m,k,n,q)$ counts weighted microstates
  --- the distinct ways $n$ spins can collectively produce a given
  pair of observables. The integrality is automatic because we are
  counting configurations of a finite system.
\end{itemize}
\end{observation}
\end{tcolorbox}

%% -------------------------------------------------------------------
\subsection{From spins to fields}

\begin{tcolorbox}[breakable, colback=yellow!10, colframe=yellow!10,
  boxrule=0pt, sharp corners]
The passage from the discrete to the continuous matrix is the passage
from statistical mechanics to quantum field theory:
\[
  \underbrace{I(m,k,n,q)}_{\text{$n$ qubits}}
  \;\xrightarrow{\;n \to \infty\;}
  \underbrace{D_{mk}}_{\text{QHO / QFT}}.
\]
In this limit:
\begin{itemize}
\item $\mathrm{Bin}(n,q) \to N(0, \sigma^2)$
  \hfill (central limit theorem)
\item Krawtchouk polynomials $\to$ Hermite polynomials
\item multinomial coefficients $\to$ matchings (Wick contractions)
\item lattice model $\to$ free field theory
\item $[1 + \alpha(s{+}u) + su]^n
  \to \exp(-(s^2 - 4su + u^2)/3)$
  \hfill ($\exp = \lim(1+x/n)^n$)
\end{itemize}

The matching matrix is a Rosetta stone between these two descriptions.
Every Feynman diagram is a limit of multinomial microstates.
\end{tcolorbox}

%% -------------------------------------------------------------------
\subsection{Quantum error correction}

\begin{tcolorbox}[colback=black!5, colframe=black!5, boxrule=0pt,
  sharp corners]
\begin{observation}[MacWilliams identity]\label{obs:mac}
The complement symmetry
$I(m,k,n,q) = \alpha^n\, I(n{-}m,\, n{-}k,\, n,\, q)$
is the MacWilliams identity of coding theory. For a quantum
error-correcting code of length $n$, the Krawtchouk transform
relates the weight distribution of the code to that of its dual.
The matching matrix generalises this: it measures how a code's
weight distribution transforms under a change of bias parameter.
\end{observation}
\end{tcolorbox}

%% ===================================================================
\section{Open Questions}
%% ===================================================================

\begin{tcolorbox}[colback=blue!5, colframe=blue!5, boxrule=0pt,
  sharp corners]
\begin{conjecture}[Discrete combinatorial interpretation]
\label{conj:discrete-comb}
For integer $\alpha \ge 1$, the entry $I(m,k,n,q)$ counts something
beyond multinomials: weighted lattice paths, coloured set partitions,
or a discrete matching analogue. A bijective proof would illuminate
the convergence to bipartite matchings as $n \to \infty$.
\end{conjecture}
\end{tcolorbox}

\begin{tcolorbox}[colback=blue!5, colframe=blue!5, boxrule=0pt,
  sharp corners]
\begin{conjecture}[Discrete diagonal generating function]
\label{conj:discrete-diag}
For fixed $\alpha$, the diagonal sequence $I(k,k,n,q)$ should have a
closed-form generating function generalising the Hermite diagonal's
$1/\sqrt{(1{-}2x)(1{-}6x)}$.
\end{conjecture}
\end{tcolorbox}

\begin{tcolorbox}[colback=blue!5, colframe=blue!5, boxrule=0pt,
  sharp corners]
\begin{conjecture}[Scaling limit]\label{conj:limit}
Under the substitutions $m = \lfloor \mu\sqrt{n}\,\rfloor$,
$k = \lfloor \kappa\sqrt{n}\,\rfloor$,
$\alpha = \beta/\sqrt{n}$, the discrete matrix converges to the
continuous:
\[
  \lim_{n \to \infty} I\!\left(\lfloor\mu\sqrt{n}\,\rfloor,\,
  \lfloor\kappa\sqrt{n}\,\rfloor,\, n,\, q\right)
  \sim D_{\mu,\kappa}^{\text{Hermite}}
  \cdot (\text{normalisation}).
\]
\end{conjecture}
\end{tcolorbox}

\begin{tcolorbox}[colback=blue!5, colframe=blue!5, boxrule=0pt,
  sharp corners]
\begin{conjecture}[Hermite matching conjecture, restated]
\label{conj:wick}
A proof of Conjecture~\ref{conj:hermite} via Isserlis' theorem
(Wick's theorem at non-unit variance) would establish that
$D_{mk}$ is a sum over Feynman diagrams --- not merely an analogy
but an identity.
\end{conjecture}
\end{tcolorbox}

%% ===================================================================
\section{Proofs Available Upon Request}
%% ===================================================================

All entries in the continuous tables were computed from the proved
coefficient-extraction formula (Theorem~\ref{thm:hermite-int}) and
independently verified against the matching
formula~(Conjecture~\ref{conj:hermite}) for all $m, k \le 20$
(441~entries, zero discrepancies). All discrete tables were computed
from the multinomial formula (Theorem~\ref{thm:kraw-coeff}) and
independently verified against the direct Krawtchouk
sum~\eqref{eq:Imk} for $n \le 12$ and all $\alpha \in \{-3, -2,
-1, 0, 1, 2, 3, \tfrac13, \tfrac12, \tfrac23, \tfrac32\}$.

\end{document}
